% Bagian: computability
% Bab: computability-theory
% Subbagian: introduction

\documentclass[../../../include/open-logic-section]{subfiles}

\begin{document}

\olfileid[id]{cmp}{thy}{int}
\olsection{Pendahuluan}

Cabang logika yang dikenal sebagai \emph{teori komputabilitas} membahas
persoalan-persoalan yang berkaitan dengan komputabilitas, atau komputabilitas
relatif, dari fungsi dan himpunan. Pengaruh Kleene tampak dari kenyataan bahwa
bidang ini dahulu dikenal sebagai \emph{teori rekursi}; dewasa ini kedua nama
tersebut sama-sama lazim digunakan.

Marilah kita menyebut suatu fungsi~$f\colon \Nat \pto \Nat$ \emph{dapat
  dihitung parsial} jika fungsi itu dapat dihitung dalam suatu model komputasi.
Jika $f$ total, kita cukup mengatakan bahwa $f$~\emph{dapat dihitung}. Suatu
relasi~$R$ yang fungsi karakteristiknya~$\Char{R}$ dapat dihitung juga disebut
dapat dihitung. Jika $f$ dan~$g$ merupakan fungsi parsial, kita menulis
$f(x) \fdefined$ untuk menyatakan bahwa $f$ terdefinisi pada~$x$, yakni bahwa
$x$~berada dalam domain~$f$; dan $f(x) \fundefined$ untuk menyatakan
kebalikannya, yakni bahwa $f$ tidak terdefinisi pada~$x$. Kita menggunakan
$f(x) \simeq g(x)$ untuk menyatakan bahwa $f(x)$ dan $g(x)$ sama-sama tidak
terdefinisi, atau keduanya terdefinisi dan bernilai sama.

Teori komputabilitas dapat dipelajari tanpa harus mengacu pada suatu model
komputasi tertentu. Untuk itu, kita menunjukkan bahwa terdapat fungsi dapat
dihitung parsial universal~$\fn{Un}(k, x)$. Fungsi ini memungkinkan kita
mengenumerasi fungsi-fungsi dapat dihitung parsial. Kita menggunakan
notasi~$\cfind{k}$ untuk menyatakan fungsi dapat dihitung parsial ke-$k$ yang
berargumen satu, yang didefinisikan oleh
$\cfind{k}(x) \simeq \fn{Un}(k, x)$.
(Kleene menggunakan $\{ k \}$ untuk tujuan ini, tetapi belakangan ini notasi
tersebut tidak banyak digunakan.) Secara sedikit lebih umum, kita dapat secara
seragam mengenumerasi fungsi-fungsi dapat dihitung parsial dengan aritas
sembarang, dan kita menggunakan $\cfind{k}[n]$ untuk menyatakan fungsi rekursif
parsial ke-$k$ yang berargumen~$n$.

Jika $f(\vec x, y)$ merupakan fungsi total atau parsial, maka $\umin{y}{f
(\vec x, y)}$ adalah fungsi dari~$\vec x$ yang mengembalikan~$y$ terkecil
sedemikian sehingga $f(\vec x, y) = 0$, dengan syarat semua nilai
$f(\vec x, 0)$, \dots, $f(\vec x, y-1)$ terdefinisi; jika tidak ada~$y$
semacam itu, $\umin{y}{f (\vec x, y)}$ tidak terdefinisi. Jika
$R(\vec x, y)$ merupakan suatu relasi, $\umin{y}{R(\vec x, y)}$ didefinisikan
sebagai~$y$ terkecil sedemikian sehingga $R(\vec x, y)$ benar; dengan kata
lain, sebagai~$y$ terkecil sedemikian sehingga
$1 \tsub \Char{R}(\vec x, y) = 0$.

Untuk menunjukkan bahwa suatu fungsi dapat dihitung, terdapat dua cara yang
dapat ditempuh:
\begin{enumerate}
\item Secara ketat: deskripsikan secara eksplisit suatu mesin Turing atau
  fungsi rekursif parsial, lalu tunjukkan bahwa mesin atau fungsi tersebut
  menghitung fungsi yang dimaksud;
\item Secara informal: deskripsikan suatu algoritma yang menghitung fungsi
  tersebut, lalu gunakan tesis Church.
\end{enumerate}
Tidak ada batas tegas di antara kedua cara itu; deskripsi algoritma yang
terperinci seharusnya memberikan informasi yang cukup sehingga relatif jelas
bagaimana, setidaknya pada prinsipnya, kita dapat merancang mesin Turing atau
rangkaian definisi rekursif parsial yang tepat. Definisi yang sepenuhnya ketat
biasanya kurang memberi pemahaman, sehingga kita akan mencoba mengambil jalan
tengah di antara kedua pendekatan tersebut. Singkatnya, kita akan mencari
bukti yang intuitif tetapi tetap ketat, sedemikian sehingga definisi yang
presisi dapat disusun darinya.

\end{document}
